\begin{name}
	{\tenchude}
	{TOÁN 12}
	{LỚP TOÁN THẦY PHÁT}
	{Thời gian: 90 phút - Không kể thời gian phát đề}
\end{name}
\TN
\Opensolutionfile{ans}[ans/ansDe2-TN1]
\begin{ex}%[EX-Ôn tập TN 2025, Nguyễn Tiến]%[2D4N1-1]
	Phát biểu nào sau đây là đúng?
	\choice
	{\True $\displaystyle\int\limits F'(x)\mathrm{\,d}x=F(x)+C$}
	{$\displaystyle\int\limits F(x)\mathrm{\,d}x=F'(x)+C$}
	{$\displaystyle\int\limits F(x)\mathrm{\,d}x=F(x)+C$}
	{$\displaystyle\int\limits F'(x)\mathrm{\,d}x=F'(x)+C$}
	\loigiai{
		Phát biểu đúng là $\displaystyle\int\limits F'(x)\mathrm{\,d}x=F(x)+C$.
	}
\end{ex}

\begin{ex}%[Đề 12 - 2025 - Nguyen Son]%[2D4H1-1]
	Cho hàm số $ f(x) $ có đạo hàm $ f'(x) $ liên tục và có một nguyên hàm là hàm số $ F(x) $. Tìm nguyên hàm $ I=\displaystyle\int \left[2f(x)+f'(x)+1\right]\mathrm{d}x $.
	\choice
	{\True $ I=2F(x)+f(x)+x+C $}
	{$ I=2F(x)+xf(x)+C $}
	{$ I=2xF(x)+f(x)+x+1 $}
	{$ I=2xF(x)+f(x)+x+C $}
	\loigiai{
		Ta có $ I=\displaystyle\int \left[2f(x)+f'(x)+1\right]\mathrm{d}x=2F(x)+f(x)+x+C $.
	}
\end{ex}

\begin{ex}%[2D4N1-2]
	Nguyên hàm của hàm số $f(x)=x^3+x$ là
	\choice
	{\True $\dfrac{1}{4}x^4+\dfrac{1}{2}x^2+C$}
	{$3x^2+1+C$}
	{$x^3+x+C$}
	{$x^4+x^2+C$}
	\loigiai{
		$\displaystyle\int{(x^3+x^2)\mathrm{\,d}x}=\dfrac{1}{4}x^4+\dfrac{1}{2}x^2+C$.}
\end{ex}

\begin{ex}%[2D4N1-3]%[Dự án EX-TF-TLN lần 3 -Mui Doan]
	Họ tất cả các nguyên hàm của hàm số $f(x)=x+\sin x$ là
	\choice
	{$\dfrac{x^2}{2}+\cos x+C$}
	{$x^2+\cos x+C$}
	{$x^2-\cos x+C$}
	{\True $\dfrac{x^2}{2}-\cos x+C$}
	\loigiai{
		Ta có $\displaystyle\int f(x)\mathrm{\,d}x=\displaystyle\int\left(x+\sin x\right)\mathrm{\,d}x=\dfrac{x^2}{2}-\cos x+C$.}
\end{ex}

\begin{ex}%[Mui Doan, dự án 12EX-OTTN2025]%[2D4N1-4]
	Hàm số nào sau đây là một nguyên hàm của hàm số $y=10^x$?
	\choice
	{$y=10^x \ln 10$}
	{$y=10^x$}
	{$y=\dfrac{10^{x+1}}{x+1}$}
	{\True $y=\dfrac{10^x}{\ln 10}$}
	\loigiai{
		Theo công thức nguyên hàm, một nguyên hàm của $y=10^x$ là $y=\dfrac{10^x}{\ln 10}$.
	}
\end{ex}

\begin{ex}%[2D4N1-4]
	$\displaystyle\int (3^x+4^x)\mathrm{d}\,x$ bằng
	\choice
	{\True  $\dfrac{3^x}{\ln 3}+\dfrac{4^x}{\ln 4}+C$}
	{$\dfrac{3^x}{\ln 4}+\dfrac{4^x}{\ln 3}+C$}
	{$\dfrac{4^x}{\ln 3}-\dfrac{3^x}{\ln 4}+C$}
	{$\dfrac{3^x}{\ln 3}-\dfrac{4^x}{\ln 4}+C$}
	\loigiai{
		Áp dụng công thức $\displaystyle\int a^x\,\mathrm{d}x=\frac{a^x}{\ln a}+C$.\\
		Ta có $\displaystyle\int(3^x+4^x)\mathrm{d}\,x
			=\int 3^x\mathrm{d}\,x+\int 4^x\mathrm{d}\,x=\dfrac{3^x}{\ln 3}+\dfrac{4^x}{\ln 4}+C$.}
\end{ex}

\begin{ex}%[2D4H1-2]
	Tìm nguyên hàm của hàm số $ f(x)=\dfrac{{x^4}+2}{x^2}$.
	\choice
	{$\displaystyle\int{f(x)\mathrm{d}x=}\dfrac{{x^3}}{3}-\dfrac{1}{x}+C$}
	{\True $\displaystyle\int{f(x)\mathrm{d}x=}\dfrac{{x^3}}{3}+\dfrac{2}{x}+C$}
	{$\displaystyle\int{f(x)\mathrm{d}x=}\dfrac{{x^3}}{3}+\dfrac{1}{x}+C$}
	{\True $\displaystyle\int{f(x)\mathrm{d}x=}\dfrac{{x^3}}{3}-\dfrac{2}{x}+C$}
	\loigiai{
		Ta có: $\displaystyle\int{f(x)\mathrm{d}x=}\displaystyle\int{\dfrac{{x^4}+2}{{x^2}}}\mathrm{d}x=\displaystyle\int{( {x^2}+\dfrac{2}{{x^2}} )}\mathrm{d}x=\dfrac{{x^3}}{3}-\dfrac{2}{x}+C$.
	}
\end{ex}

\begin{ex}%[2D4H1-3]
	$\displaystyle \int \left(\cos \dfrac{x}{2} \right)^2 \mathrm{d}x$ bằng
	\choice
	{$x+\sin x+C$}
	{$\dfrac{1}{3}\left(\cos \dfrac{x}{2} \right)^3+C$}
	{$\left(\sin \dfrac{x}{2} \right)^2+C$}
	{\True $\dfrac{1}{2}x+ \dfrac{1}{2}\sin x+C$}
	\loigiai{
		Ta có
		\allowdisplaybreaks
		\begin{eqnarray*}
			\displaystyle \int \left(\cos \dfrac{x}{2} \right)^2 \mathrm{d}x &=& \displaystyle \int \dfrac{1+\cos x}{2} \mathrm{d}x\\
			&=& \dfrac{1}{2} \displaystyle \int \mathrm{d}x+ \dfrac{1}{2}\displaystyle \int \cos x \mathrm{d}x \\
			&=& \dfrac{1}{2}x+ \dfrac{1}{2}\sin x+C.
		\end{eqnarray*}
	}
\end{ex}

\begin{ex}%[Dự án 2025 - đề cấu trúc mới, Nguyễn Kiều Nhã Tú]%[2D4N2-1]
	Cho hai hàm số $f(x)$ và $g(x)$ liên tục trên $K$, $a$, $b\in K$. Khẳng định nào sau đây là \textbf{sai}?
	\choice
	{$\displaystyle\int\limits_a^b [f(x)+g(x)]\mathrm{\,d} x=\displaystyle\int\limits_a^b f(x)\mathrm{\,d} x+\displaystyle\int\limits_a^b g(x)\mathrm{\,d}x$}
	{$\displaystyle\int\limits_a^b 2f(x)\mathrm{\,d}x=2 \displaystyle\int\limits_a^b f(x)\mathrm{\,d}x$}
	{\True $\displaystyle\int\limits_a^b f(x)g(x)\mathrm{\,d} x=\displaystyle\int\limits_a^b f(x)\mathrm{\,d}x \cdot \displaystyle\int\limits_a^b g(x)\mathrm{\,d}x$}
	{$\displaystyle\int\limits_a^b[f(x)-g(x)]\mathrm{\,d} x=\displaystyle\int\limits_a^b f(x)\mathrm{\,d} x-\displaystyle\int\limits_a^b g(x)\mathrm{\,d}x$}
	\loigiai{
		Theo tính chất của tích phân, không có tính chất $\displaystyle\int\limits_a^b f(x)g(x)\mathrm{\,d} x=\displaystyle\int\limits_a^b f(x)\mathrm{\,d}x \cdot \displaystyle\int\limits_a^b g(x)\mathrm{\,d}x$.
	}
\end{ex}

\begin{ex}%[2D4N2-2]%[To 20 - Dot 17 - Chuong 4 - Bai 3 - CD - De 3]%[Đình Nguyên]
	Tính $\displaystyle\int\limits_{3}^{5}\left(x^{2}+2 x\right) \mathrm{\,d}x$.
	\choice
	{\True $\dfrac{146}{3}$}
	{$146$}
	{$3$}
	{$\dfrac{143}{6}$}
	\loigiai
	{
		Ta có
		\begin{align*}
			\int\limits_{3}^{5}\left(x^{2}+2 x\right) \mathrm{\,d}x & =\int\limits_{3}^{5} x^{2} \mathrm{\,d}x+2 \int\limits_{3}^{5} x \mathrm{\,d}x             \\
			                                                        & =\left.\dfrac{1}{3} x^{3}\right|_{3} ^{5}+\left.2 \cdot \dfrac{1}{2} x^{2}\right|_{3} ^{5} \\
			                                                        & =\left.\dfrac{1}{3} x^{3}\right|_{3} ^{5}+\left.x^{2}\right|_{3} ^{5}                      \\
			                                                        & =\dfrac{1}{3}\left(5^{3}-3^{3}\right)+\left(5^{2}-3^{2}\right)=\dfrac{146}{3}.
		\end{align*}
	}
\end{ex}

\begin{ex}%[2D4N2-3]
	Tích phân $I=\displaystyle\int\limits_{\frac{\pi}{4}}^{\frac{\pi}{3}} \dfrac{\mathrm{\,d} x}{\sin ^2 x}$ bằng
	\choice
	{$\cot \dfrac{\pi}{3}-\cot \dfrac{\pi}{4}$}
	{$\cot \dfrac{\pi}{3}+\cot \dfrac{\pi}{4}$}
	{\True $-\cot \dfrac{\pi}{3}+\cot \dfrac{\pi}{4}$}
	{$-\cot \dfrac{\pi}{3}-\cot \dfrac{\pi}{4}$}
	\loigiai{Ta có
	$I=\displaystyle\int\limits_{\frac{\pi}{4}}^{\frac{\pi}{3}} \dfrac{\mathrm{\,d} x}{\sin ^2 x}=-\left.\cot x\right|_{\frac{\pi}{4}} ^{\frac{\pi}{3}}=-\cot \dfrac{\pi}{3}+\cot \dfrac{\pi}{4}$.
	}
\end{ex}

\begin{ex}%[Tổ 20 - Chương 4 - - CD]%[Nguyễn Văn Sang]%[2D4N2-4]
	Tính $I=\displaystyle\int\limits_0^1\left(2 \cdot  \mathrm{e}^{-x}+4^x\right) \mathrm{\,d}x$
	\choice
	{\True $I=\dfrac{-2}{\mathrm{e}}+2+\dfrac{3}{\ln 4}$ }
	{$I=\dfrac{2}{\mathrm{e}}+2+\dfrac{3}{\ln 4}$ }
	{$I=\dfrac{-2}{\mathrm{e}}+2-\dfrac{3}{\ln 4}$ }
	{$I=\dfrac{-2}{\mathrm{e}}+2+\dfrac{4}{\ln 4}$ }
	\loigiai{
		\[I=\displaystyle\int\limits_0^1\left(2 \cdot \mathrm{e}^{-x}+4^x\right) \mathrm{\,d} x=\left(-2 \mathrm{e}^{-x}+\dfrac{4^x}{\ln 4}\right)\bigg|_0 ^1=\dfrac{-2}{\mathrm{e}}+\dfrac{4}{\ln 4}+2-\dfrac{1}{\ln 4}=\dfrac{-2}{\mathrm{e}}+2+\dfrac{3}{\ln 4}.\]
	}
\end{ex}
\Closesolutionfile{ans}

\TNTF
\Opensolutionfile{ans}[ans/ansDe2-TN2]
\begin{ex}%[BG-12-4in1, Phạm Đức]%[2D4N1-1]
	Cho hàm số $f(x)$ có một nguyên hàm là $F(x)$, $k$ là số thực bất kỳ. Mỗi khẳng định sau đây đúng hay sai?
	\choiceTF[t]
	{\True $\displaystyle\int f(x)\mathrm{\,d}x=F(x)+C$ với $C$ là hằng số}
	{$f'(x)=F(x)$}
	{$\displaystyle\int kf(x)\mathrm{\,d}x=k\int f(x)\mathrm{\,d}x$}
	{\True $\displaystyle\int \left[f(x)+F'(x)\right]\mathrm{\,d}x=2\left[F(x)+2025\right]+C$ với $C$ là hằng số}
	\loigiai{
		\begin{itemchoice}
			\itemch Theo định nghĩa nguyên hàm ta có $\displaystyle\int f(x)\mathrm{\,d}x=F(x)+C$ với $C$ là hằng số.
			\itemch Tổng quát ta chỉ có kết quả $F'(x)=f(x)$.
			\itemch Tính chất $\displaystyle\int kf(x)\mathrm{\,d}x=k\int f(x)\mathrm{\,d}x$ chỉ đúng với $k\ne 0$.
			\itemch Ta có $\left(2\left[F(x)+2025\right]+C\right)'=2F'(x)=2f(x)=f(x)+f(x)=f(x)+F'(x)$.

			Vậy $\displaystyle\int \left[f(x)+F'(x)\right]\mathrm{\,d}x=2\left[F(x)+2025\right]+C$ với $C$ là hằng số.
		\end{itemchoice}
	}
\end{ex}

\begin{ex}%[2D4H2-1]
	Cho $\displaystyle\int\limits_0^{\frac{\pi}{2}} \left[\sin x +f(x)\right] \mathrm{\, d}x=6$. Xét tính đúng, sai cho mỗi khẳng định sau.
	\choiceTF[t]
	{\True $\displaystyle\int\limits_0^{\tfrac{\pi}{2}} \left[\sin x +f(x)\right] \mathrm{\, d}x= \displaystyle\int\limits_0^{\tfrac{\pi}{2}} \sin x \mathrm{\, d}x + \displaystyle\int\limits_0^{\tfrac{\pi}{2}} f(x) \mathrm{\, d}x$}
	{\True  $\displaystyle\int\limits_0^{\tfrac{\pi}{2}} \sin x\mathrm{\,d}x=1$}
	{\True $\displaystyle\int\limits_0^{\tfrac{\pi}{2}} f(x)\mathrm{\, d}x=5$}
	{$\displaystyle\int\limits_0^{\tfrac{\pi}{2}}\left[3+2f(x)\right]\mathrm{\, d}x=\dfrac{3\pi}{2}+5$}
	\loigiai{
		\begin{itemchoice}
			\itemch  Ta có $\displaystyle\int\limits_0^{\tfrac{\pi}{2}} \left[\sin x +f(x)\right] \mathrm{\, d}x= \displaystyle\int\limits_0^{\tfrac{\pi}{2}} \sin x \mathrm{\, d}x + \displaystyle\int\limits_0^{\tfrac{\pi}{2}} f(x) \mathrm{\, d}x$.
			\itemch  Ta có $\displaystyle\int\limits_0^{\tfrac{\pi}{2}} \sin x\mathrm{\,d}x=-\cos x\bigg|_0^{\tfrac{\pi}{2}}=-\cos \dfrac{\pi}{2}+\cos 0=1$.
			\itemch  Ta có $\displaystyle\int\limits_0^{\tfrac{\pi}{2}} \left[\sin x +f(x)\right] \mathrm{\, d}x= \displaystyle\int\limits_0^{\tfrac{\pi}{2}} \sin x \mathrm{\, d}x + \displaystyle\int\limits_0^{\tfrac{\pi}{2}} f(x) \mathrm{\, d}x= 1 + \displaystyle\int\limits_0^{\tfrac{\pi}{2}} f(x) \mathrm{\, d}x=6$.\\
			Suy ra $\displaystyle\int\limits_0^{\tfrac{\pi}{2}} f(x)\mathrm{\, d}x=6-1=5$.
			\itemch  Ta có $\displaystyle\int\limits_0^{\tfrac{\pi}{2}}\left[3+2f(x)\right]\mathrm{\, d}x=\displaystyle\int\limits_0^{\tfrac{\pi}{2}}3\mathrm{\, d}x+2\displaystyle\int\limits_0^{\tfrac{\pi}{2}} f(x)\mathrm{\, d}x=3x\bigg|_0^{\tfrac{\pi}{2}}+2\cdot5 =\dfrac{3\pi}{2}+10\neq\dfrac{3\pi}{2}+5$.
		\end{itemchoice}
	}
\end{ex}
\Closesolutionfile{ans}

\TNSA
\Opensolutionfile{ans}[ans/ansDe2-TN3]
\begin{ex}%[2D4H1-1]%[Đào Trung Kiên]
	Biết $F(x)$ là một nguyên hàm của hàm số $f(x)=\sin x$ và đồ thị hàm số $y=F(x)$ đi qua điểm $M\left(0;1\right)$. Tính $F\left(\dfrac{\pi}{2}\right)$ (làm tròn kết quả tới hàng đơn vị).
	\shortans[]{$2$}
	\loigiai{
		Ta có $F(x)=\displaystyle\int f(x)\mathrm{\,d}x=-\cos x+C$.\\
		Mà đồ thị hàm số $y=F(x)$ đi qua $M(0;1)$ nên $F(0)=1\Leftrightarrow -1+C=1\Leftrightarrow C=2$.\\
		Suy ra $F(x)=-\cos x+2$ nên $F\left(\dfrac{\pi}{2}\right)=2$.}
\end{ex}

\begin{ex}%[2D4H2-2]%[Tổ 20 - Đợt 17 - Chương 4 - - CD]%[Phạm Hà Giang]
	Tích phân $\displaystyle\int\limits_2^8 \dfrac{1}{x} \mathrm{\,d}x=a\ln 2$. Giá trị của $a$ bằng bao nhiêu?
	\shortans{$2$}
	\loigiai
	{ $\displaystyle\int\limits_2^8 \dfrac{1}{x} \mathrm{\,d}x =\left. \ln \left| x \right| \right| _2^8=\ln 8 -\ln 2 =2\ln 2.$\\
		Suy ra $a=2$.
	}
\end{ex}

\begin{ex}%[2D4H3-1]
	\immini{Cho hàm số $f(x)$ liên tục trên $\mathbb{R}$. Đồ thị hàm số $y=f'(x)$ được cho như hình bên. Diện tích các hình phẳng $(K)$, $(H)$ lần lượt là $\dfrac{5}{12}$ và $\dfrac{8}{3}$. Biết $f(-1)=\dfrac{19}{12}$, tính $f(2)$(kết quả làm tròn đến hàng phần mười).
	}
	{
		\begin{tikzpicture}[scale=1, font=\footnotesize, line join=round, line cap=round, >=stealth]
			\tikzset{label style/.style={font=\footnotesize}}
			%Nhập giới hạn đồ thị và hàm số cần vẽ
			\def \xmin{-1.5}
			\def \xmax{2.5}
			\def \ymin{-2.6}
			\def \ymax{1}
			\def \hamso{sin((2*(\x))*180/pi)}
			%\def \tiemcanxien{\x+1}
			%Tự động
			\draw[->] (\xmin,0)--(\xmax,0) node[below left] {$x$};
			\draw[->] (0,\ymin)--(0,\ymax) node[below left] {$y$};
			\draw[fill=black] (0,0) circle(1pt) node [below left] {$O$};
			%Vẽ các điểm trên 2 hệ trục
			\foreach \x in {-1,2}
			\fill[black] (\x,0) circle(1pt) node [above left] {$\x$};
			%Tự động
			\begin{scope}
				\clip (\xmin+0.01,\ymin+0.01) rectangle (\xmax-0.01,\ymax-0.01);
				\draw[samples=350,domain=\xmin+0.01:\xmax-0.01,smooth,variable=\x] plot (\x,{(\x+1)*(\x)*(\x-2)});
			\end{scope}
			\draw[pattern = north east lines] (-1,0)--plot[domain=-1:2] (\x,{(\x+1)*(\x)*(\x-2)});
			\draw (-.5,-.5)node{$(K)$};
			\draw (1,.5)node{$(H)$};
		\end{tikzpicture}
	}
	\shortans[]{$-0{,}7$}
	\loigiai{
	Ta có $S_K=\displaystyle \int\limits_{-1}^0 f'(x)\mathrm{\,d}x
		=f(0)-f(-1)	=\dfrac{5}{12}
		\Rightarrow f(0)=\dfrac{5}{12}+f(-1)=\dfrac{5}{12}+\dfrac{19}{12}=2$.\\
	Lại có: $S_H=\displaystyle \int\limits_0^2 \left| f'(x) \right|\mathrm{\,d}x
		=-\displaystyle \int\limits_0^2 f'(x)\mathrm{\,d}x
		=f(0)-f(2)=\dfrac{8}{3}
		\Rightarrow f(2)=f(0)-\dfrac{8}{3}=2-\dfrac{8}{3}=\dfrac{-2}{3}\approx -0{,}7$.
	}
\end{ex}

\begin{ex} %[2D4H3-3]
	Cho hình phẳng $(D)$ được giới hạn bởi các đường $ y={x^2}+2x+2$; $y=6-x$; $y=2$ và $(D)$ nằm ngoài Parabol $y={x^2}+2x+2$. Khi cho $(D)$ quay quanh trục $Ox$, ta nhận được vật thể tròn xoay có thể tích $V=\dfrac{a\pi }{b}$, trong đó $a$, $b$ là các số nguyên dương. Giá trị biểu thức $P=a-2{b^2}$ bằng bao nhiêu.
	\shortans[1]{$73$}
	\loigiai{
	Vẽ các đường $ y={x^2}+2x+2$; $y=6-x$; $y=2$.
	\begin{center}
		\begin{tikzpicture}[scale=0.7,line join=round,line cap=round,font=\footnotesize,>=stealth]
			\def\f(#1){1*((#1)^2)+2*(#1)+2}
			\def\g(#1){6-1*(#1)}
			\draw[->] (-3,0)--(7,0) node[below] { $x$};
			\draw[->] (0,-1)--(0,7) node[left] { $y$};
			\draw (0,0) node [below left] { $O$};
			\foreach \x in {1,4}		\draw (\x,0.1)--(\x,-0.1) node [below] { $\x$};
			\foreach \y in {5}		\draw (0.1,\y)--(-0.1,\y) node [left] { $\y$};
			\node[left] at (0,2.2) {$2$};
			\clip (-4,-1) rectangle (6,6);
			\draw[thick,smooth,samples=200] plot[domain=-3:3] (\x,{\f(\x)});
			\draw[thick,smooth,samples=200] plot[domain=-2:7] (\x,{\g(\x)});
			\draw[thick,smooth,samples=200] plot[domain=-3:3] (-3,2)--(6,2);
			\draw[dashed] (4,0)--(4,2) (1,0)-- (1,5)--(0,5);
			\fill[pattern=north east lines,pattern color=red]plot[domain=0:2](\x,{(\x)^2+2*(\x)+2})-- (1,5)--(1,2) -- cycle;
			\fill[pattern=north east lines,pattern color=black] (1,2)--(1,5)--plot[domain=0:2](\x,{6-1*(\x)})-- (4,2) -- cycle;
		\end{tikzpicture}
	\end{center}
	Dựa vào hình vẽ ta có\\
	$V=\pi \int\limits_0^1{| {{( {x^2}+2x+2 )}^2}-{2^2} |\mathrm{\,d}x}+\pi \int\limits_1^4{| {{( 6-x )}^2}-{2^2} |\mathrm{\,d}x}=\dfrac{523\pi }{15}\\
		\Rightarrow a=523,b=15$.\\
	Vậy $P=523-{{2\cdot 15}^2}=73$.
	}
\end{ex}

\Closesolutionfile{ans}

\TL
\begin{ex}%[2D4H2-2] 
	Tính tích phân $I=\displaystyle\int\limits_{-1}^2{\left|x^3-3x+2\right|}\mathrm{\,d}x$.
	\loigiai{
	$C=\displaystyle\int\limits_{-1}^2 {\left|x^3-3x+2\right|}\mathrm{\,d}x$.\\
	Xét $f(x) = x^3-3x-2$ trên $[-1; 2]$.\\
	Cho $f(x) = 0
		\Leftrightarrow
		x^3-3x-2 = 0
		\Leftrightarrow
		\hoac{& x=2\in [-1; 2] \\ & x=-1\in [-1; 2].}$\\
	Bảng xét dấu $f(x)$ trên đoạn $[-1; 2]$.
	\begin{center}
		\begin{center}
			\begin{tikzpicture}
				\tkzTabInit[nocadre=false,lgt=3,espcl=2.5,deltacl=0]
				{$x$ /0.6,$x^3-3x+2$ /0.6}
				{,$-1$,$2$,}
				\tkzTabLine{,h,$0$,-,$0$,h,}
			\end{tikzpicture}
		\end{center}
	\end{center}
	Do đó
	$C = -\displaystyle\int\limits_{-1}^2 {(x^3-3x-2)}\mathrm{\,d}x	=
		-\left.\left(\dfrac{1}{4} x^4-\dfrac{3}{2} x^2 -2x\right)\right|_{-1}^2=
		6- \left( -\dfrac{3}{4} \right)=\dfrac{27}{4}$.
	}
\end{ex}


\begin{ex}%[2D4V1-4]
	Cho hàm số $f(x)$ thỏa mãn $ f(0)=1-\ln 2$ và $\mathrm{e}^ xf'(x)=2^x\left[f(x)\right]^2$ với mọi $x\in\mathbb{R}$. Giá trị của $f(1)$ bằng bao nhiêu? (\textit{Kết quả làm tròn đến hàng phần trăm}).
	% \shortans{$0{,}42$}
	\loigiai{
		Từ giả thiết ta có $f'(x)=\dfrac{2^x}{\mathrm{e}^ x}{\left[f(x)\right]^2}$ với mọi $ x\in\left(1;2\right]$.\\
		Do đó $ f(x)\ge f(1)=1>0$ với mọi $ x\in\left[1;2\right]$.\\
		Xét với mọi $ x\in [1 ; 2]$ ta có
		\begin{align*}
			\mathrm{e}^ x{f}'(x)=2^x{\left[f(x)\right]^2} & \Rightarrow{f}'(x)=\dfrac{2^x}{\mathrm{e}^ x}{\left[f(x)\right]^2} \\&\Rightarrow\dfrac{f'(x)}{\left[f(x)\right]^2}=\left(\dfrac{2}{\mathrm{e}}\right)^x\\&\Rightarrow-\left(\dfrac{1}{f(x)}\right)'=\left(\dfrac{2}{\mathrm{e}}\right)^x \\&\Rightarrow{\left(\dfrac{1}{f(x)}\right)'}=-\left(\dfrac{2}{\mathrm{e}}\right)^x\\&\Rightarrow\dfrac{1}{f(x)}=-\displaystyle\int\left(\dfrac{2}{\mathrm{e}}\right)^x\mathrm{\,d} x\\&\Rightarrow\dfrac{1}{f(x)}=-\dfrac{\left(\dfrac{2}{\mathrm{e}}\right)^x}{\ln \dfrac{2}{\mathrm{e}}}+C\\&\Rightarrow\dfrac{1}{f(x)}=\dfrac{\left(\dfrac{2}{\mathrm{e}}\right)^x}{1-\ln 2}+C.
		\end{align*}
		Mà $ f(0)=1-\ln 2\Rightarrow C=0$. \\Do đó
		$\dfrac{1}{f(x)}=\dfrac{\left(\dfrac{2}{\mathrm{e}}\right)^x}{1-\ln 2}$
		$\Rightarrow f(x)=\dfrac{1-\ln 2}{\left(\dfrac{2}{\mathrm{e}}\right)^x}=\dfrac{(1-\ln 2)\mathrm{e}^x}{2^x}$.\\
		Vậy $ f(1)=\dfrac{\mathrm{e}-\mathrm{e}\ln 2}{2}\approx 0{,}42$.}
\end{ex}

\begin{ex}%[12-MH-2-MH2025]%[MH-2025, Nguyễn Trần Phong]%[2D4C3-2]
	\immini{Chướng ngại vật \lq\lq  tường cong\rq\rq trong một sân thi đấu X-Game là một khối bê tông có chiều cao từ mặt đất lên là $3$ m. Giao của mặt tường cong và mặt đất là đoạn thẳng $AB = 2$ m. Thiết diện của khối tường cong cắt bởi mặt phẳng vuông góc với $AB$ tại $A$ là một hình tam giác vuông cong $ACE$ với $AC = 4$ m, $CE = 3$ m và cạnh cong $AE$ nằm trên một đường Parabol có trục đối xứng vuông góc với mặt đất. Tại vị trí $M$ là trung điểm của $AC$ thì tường cong có độ cao $1$ m. Thể tích bê tông cần sử dụng để tạo nên khối tường cong đó gần nhất với số nào dưới đây?
		% \shortans{$9{,}3$}
	}{\begin{tikzpicture}[>=stealth,x=0.8cm,y=0.8cm,scale=0.7]
			\coordinate[label=below:$A$] (A) at (0,0);
			\coordinate[label=left:$B$] (B) at (-2,2);
			\coordinate[label=below:$C$] (C) at (6,0);
			\coordinate[label=right:$E$] (E) at (6,6);
			\coordinate (G) at (2,4);
			\coordinate (H) at (4,1.8);
			\coordinate (D) at ($(C)+(B)-(A)$);
			\coordinate (F) at ($(E)+(D)-(C)$);
			\coordinate[label=below:$M$] (M) at ($(A)!0.5!(C)$);
			\coordinate (K) at ($(A)!0.5!(B)$);
			\coordinate (N) at ($(M)+(0,1.1)$);
			\draw (D)--(C)--(A)--(B) (C)--(E)--(F) (M)--(N);
			\draw[dashed] (B)--(D)--(F);
			\foreach \diem in {A,B,C,D,E,F,M,F}	\fill (\diem)circle(1.5pt);
			%\tkzLabelPoints[above left](D)
			%\tkzLabelSegment[right](M,N){\footnotesize$1$ m}
			%\tkzLabelSegment[left](A,B){\footnotesize$2$ m}
			%\tkzLabelSegment[right](C,E){\footnotesize$3{,}5$ m}
			\draw(-1,.8) node[left]{\footnotesize $2$ m} (3,0.8) node[right]{\footnotesize$1$ m} (6,3) node[right]{\footnotesize$3$ m};

			\draw plot[smooth,tension=.65] coordinates{(B) (G) (F)};
			\draw plot[smooth,tension=.65] coordinates{(A) (H) (E)};
			\fill [pattern = north east lines] plot[smooth,tension=.65] coordinates{(A) (H) (E)} (0,0) --(-2,2)--(4,8)--(6,6)--cycle;
			\fill [draw=none, pattern = north east lines, color=white] (0,0) plot[smooth,tension=.65] coordinates{(B) (G) (F)} (-2,2)--(4,2)--cycle;
		\end{tikzpicture}
	}
	\loigiai{
		\immini{Chọn hệ trục tọa độ như hình vẽ.\\
			Gọi $AE \colon y = ax^2 + bx + c$.\\
			Do $AE$ đi qua $A(-4; 0)$ nên ta có $16a - 4b + c = 0$.\\
			Do $E (0; 3)$ thuộc cạnh cong $AE$ nên $c = 3$ (2).\\
			Do $N(-2; 1)$ thuộc cạnh cong $AE$ nên $4a - 2b + c = 1$ (3).\\
			Từ (1), (2), (3) suy ra $a = \dfrac{1}{8}$, $b = \dfrac{5}{4}$, $c = 3 \Rightarrow AE \colon y = \dfrac{1}{8}x^2 + \dfrac{5}{4}x + 3$.\\
			Khi đó $S_{AEC} = \displaystyle\int_{-4}^0\left(\dfrac{1}{8} x^2 + \dfrac{5}{4}x + 3\right) dx = \dfrac{14}{3}\left(m^2\right)$.
		}{\begin{tikzpicture}[>=stealth,x=0.8cm,y=0.8cm,scale=0.7]
				\coordinate[label=below:$A$] (A) at (0,0);
				\coordinate[label=left:$B$] (B) at (-2,2);
				\coordinate[label=below:$C$] (C) at (6,0);
				\coordinate[label=right:$E$] (E) at (6,6);
				\coordinate (G) at (2,4);
				\coordinate (H) at (4,1.8);
				\coordinate[label = above left:$D$] (D) at ($(C)+(B)-(A)$);
				\coordinate[label = above:$F$] (F) at ($(E)+(D)-(C)$);
				\coordinate[label=below:$M$] (M) at ($(A)!0.5!(C)$);
				\coordinate (K) at ($(A)!0.5!(B)$);
				\coordinate[label=above left:$N$] (N) at ($(M)+(0,1.1)$);
				\draw (D)--(C)--(A)--(B) (C)--(E)--(F) (M)--(N);
				\draw[dashed] (B)--(D) (D)--(F);
				\foreach \diem in {A,B,C,D,E,F,M,F,N}	\fill (\diem)circle(1.5pt);
				\coordinate (x) at ($(M)!1.5!(C)$);
				\draw[->](C)--(x); \draw (x) node[right]{$x$};
				\coordinate (y) at ($(C)!1.4!(E)$);
				\draw[->](E)--(y); \draw (y) node[right]{$y$};
				%\tkzLabelPoints[above left](D)
				%\tkzLabelSegment[right](M,N){\footnotesize$1$ m}
				%\tkzLabelSegment[left](A,B){\footnotesize$2$ m}
				%\tkzLabelSegment[right](C,E){\footnotesize$3{,}5$ m}
				\draw plot[smooth,tension=.65] coordinates{(B) (G) (F)};
				\draw plot[smooth,tension=.65] coordinates{(A) (H) (E)};
				\fill [pattern = north east lines] plot[smooth,tension=.65] coordinates{(A) (H) (E)} (0,0) --(-2,2)--(4,8)--(6,6)--cycle;
				\fill [draw=none, pattern = north east lines, color=white] (0,0) plot[smooth,tension=.65] coordinates{(B) (G) (F)} (-2,2)--(4,2)--cycle;
			\end{tikzpicture}
		}
		\noindent Thể tích khối tường cong là $V = S_{AEC} \cdot AB = \frac{14}{3} \cdot 2 = \dfrac{28}{3} = 9{,}3\left(\mathrm{~m}^3\right)$.	}
\end{ex}

% \Closesolutionfile{ansbook}
% \HetDe
% \label{De2}
% %
% \cleardoublepage
% \setcounter{page}{1}
% \rfoot{Trang \thepage/\pageref{DA2} - Đáp án trắc nghiệm Mã đề 2}
% \begin{center}
% 	\bfseries ĐÁP ÁN TRẮC NGHIỆM MÃ ĐỀ 2
% \end{center}

% \inputansbox{10}{ans/ansDe2-TN1}
% \inputansbox[3]{2}{ans/ansDe2-TN2}
% \inputansbox{3}{ans/ansDe2-TN3}
% \label{DA2}
% %
